% Generated draft: requires mathematical and source review.
\documentclass[11pt]{article}
\usepackage[margin=25mm]{geometry}
\usepackage{fontspec}
\setmainfont{Latin Modern Roman}
\usepackage{amsmath,amssymb,mathtools,unicode-math}
\setmathfont{Latin Modern Math}
\usepackage{enumitem,needspace,xurl,hyperref}
\setlength{\parindent}{0pt}
\setlength{\parskip}{5pt}
\setlength{\emergencystretch}{3em}
\newcommand{\N}{\mathbb{N}}
\newcommand{\Z}{\mathbb{Z}}
\newcommand{\Q}{\mathbb{Q}}
\newcommand{\R}{\mathbb{R}}
\newcommand{\C}{\mathbb{C}}
\newcommand{\F}{\mathbb{F}}
\newcommand{\A}{\mathbb{A}}
\newcommand{\PP}{\mathbb{P}}
\newcommand{\E}{\mathbb{E}}
\newcommand{\eps}{\varepsilon}
\newcommand{\dd}{\,\mathrm{d}}
\begin{document}
\section*{UnsolvedMath ID: 10600011}
\textbf{Source problem code:} catalogue entry
\subsection{Definitions and mathematical statement}
Let \(\sim_v\) be virtual-knot equivalence, generated by virtual Reidemeister and detour moves, and let \(\sim_c\) denote classical knot equivalence. A virtual knot is classical if it has a representative with no virtual crossings. Let \(V\) be the Jones polynomial normalized by \(V(U)=1\) for the unknot. In this item, crossing virtualization is explicitly described as flanking a classical crossing by two virtual crossings; the source states that this operation preserves \(V\). Let \(\approx\) be the equivalence relation generated by \(\sim_v\) and that operation.

The compound item asks:
\begin{enumerate}
\item Can there exist a nontrivial classical knot diagram \(K\) and a subset \(S\) of its crossings such that virtualizing those crossings produces a virtual knot \(\widetilde K\) with
\[
 \widetilde K\not\sim_v U,\qquad V(\widetilde K)=1,
 \qquad \widetilde K\sim_v C
\]
for some classical knot diagram \(C\)? The starting knot must also be nontrivial, as printed.
\item For every virtual knot diagram \(K\) with \(V(K)=1\), is there a classical diagram \(C\) such that \(K\approx C\)?
\item For all classical knots \(K,K'\), does \(K\approx K'\) imply \(K\sim_c K'\)?
\end{enumerate}
The source observes that affirmative answers to the last two questions would imply that the only classical knot with unit Jones polynomial is the unknot. This is a supplied conditional consequence, not an asserted resolution.

The exact local virtualization convention must be kept consistent: the first paragraph refers simply to virtualizing a crossing, whereas the second explicitly uses the two-flanking-virtual-crossings description and states Jones invariance. A different convention involving an additional classical crossing switch is not silently substituted. Full local diagrams are absent, as are explicit orientation and polynomial-variable conventions.
\subsection{Short English statement}
Ask whether a nontrivial unit-Jones virtual knot obtained from a nontrivial classical knot by virtualization can be classical; whether every unit-Jones virtual knot becomes classical under virtualization moves; and whether those moves ever identify distinct classical knots.
\subsection{Sources}
\begin{itemize}
\item \emph{Unsolved Problems in Virtual Knot Theory and Combinatorial Knot Theory}, revised survey, problem-list item 11, ``Virtualization and Unit Jones Polynomial.'' \url{https://arxiv.org/abs/1409.2823}. Recorded access 2026-07-29; extraction source-tex. Individual authors are not supplied in this record.
\item The research summary mentions related Jones-polynomial work on long virtual knots (2020, arXiv:2012.xxxx) and checkerboard-colourable virtual knots (2000, arXiv:0008.xxxx). These are incomplete placeholder identifiers, with no complete author/title citations supplied.
\item Local UnsolvedMath registry, \texttt{data/unsolvedmath/problems.json}, permanent ID \(10600011\), exact source code \texttt{catalogue entry}; supplied OPEN-TRIAGE status.
\end{itemize}
\end{document}
